\documentclass[11pt,letterpaper]{article}
\usepackage[T1]{fontenc}
\usepackage{lmodern,microtype}
\usepackage[margin=.78in,headheight=15pt,headsep=15pt]{geometry}
\usepackage{mathtools,amssymb,booktabs,array,xcolor}
\usepackage[most]{tcolorbox}
\usepackage{enumitem,fancyhdr,needspace,hyperref}
\definecolor{ink}{HTML}{17243A}
\definecolor{indigo}{HTML}{293D70}
\definecolor{mint}{HTML}{2A8C82}
\definecolor{muted}{HTML}{647084}
\definecolor{mist}{HTML}{EAF5F2}
\color{ink}
\hypersetup{colorlinks=true,linkcolor=mint,urlcolor=mint,
 pdfauthor={Igor Rivin},pdftitle={Probability Problem Set 3: Counting, Collisions, and Derangements}}
\setlength{\parindent}{0pt}
\setlength{\parskip}{4pt}
\clubpenalty=10000
\widowpenalty=10000
\setlist[enumerate]{label=(\alph*),leftmargin=1.8em,itemsep=3pt,topsep=4pt}
\setlength{\abovedisplayskip}{7pt}
\setlength{\belowdisplayskip}{7pt}
\renewcommand{\P}{\mathbb P}
\newcommand{\pts}[1]{\hfill\textcolor{muted}{\small[#1 points]}}
\newcommand{\partpts}[1]{\unskip\nobreak\hspace{.3em}\mbox{\textcolor{muted}{\small(#1)}}}
\newcommand{\edition}{\ifsolutions Instructor solutions\else Student handout\fi}
\newcounter{problem}
\newcommand{\problem}[2]{\par\addvspace{11pt}\Needspace{7\baselineskip}\stepcounter{problem}{\bfseries\color{indigo}\theproblem.\ #1}\pts{#2}\par\nobreak}
\newcommand{\topic}[1]{\par\addvspace{9pt}\Needspace{10\baselineskip}{\large\bfseries\color{mint}#1}\par\nobreak}
\newcommand{\studentbreak}{\ifsolutions\else\newpage\fi}
\pagestyle{fancy}
\fancyhf{}
\fancyhead[L]{\small\color{muted}Probability \textbar\ Problem Set 3}
\fancyhead[R]{\small\color{muted}Chance, Data, and Decisions}
\fancyfoot[L]{\footnotesize\color{muted}\edition}
\fancyfoot[R]{\small\color{mint}\thepage}
\renewcommand{\headrulewidth}{.3pt}
\NewDocumentEnvironment{solution}{+b}{\ifsolutions
 \begin{tcolorbox}[enhanced,breakable,colback=mist!45,colframe=mint!45,
  boxrule=.5pt,arc=1mm,left=3mm,right=3mm,top=2mm,bottom=2mm,
  before skip=7pt,after skip=7pt,title={Solution},
  title after break={Solution (continued)},colbacktitle=mist,coltitle=mint,
  fonttitle=\small\bfseries]\small #1\end{tcolorbox}\fi}{}

\begin{document}
{\LARGE\bfseries\color{indigo}Problem Set 3}\hfill{\small\color{mint}\edition}\\[4pt]
{\large Counting, collisions, and derangements}\\[4pt]
{\small\color{muted}100 points + 10 optional challenge points}\par
\ifsolutions
\medskip
\textbf{Teaching note.} This assignment consolidates discrete probability
spaces and counting, then pushes the birthday and hat-check problems one
step past the lectures: rigorous two-sided birthday bounds, collision
patterns, the derangement recurrence, and the distribution of the number of
correct hats. Allow approximately three and a half hours, excluding the
challenge. Subpart scores appear in parentheses. Give credit for any correct
model and reasoning; distinguish a wrong model from an arithmetic error
carried through correctly.
\else
\medskip
\textbf{Name:}\ \rule{.46\linewidth}{.4pt}\hfill
\textbf{Due:}\ \rule{.20\linewidth}{.4pt}
\fi

\medskip
\textbf{Instructions.} Identify the outcomes, their probabilities, and the
event being counted. State whether order matters and repetition is allowed.
Give exact answers except where decimal approximations are requested.
A calculator or short program may be used for Problems 3, 6, and 7.
Throughout, $D_n$ is the number of derangements of $n$ objects
(permutations with no fixed point), and $p_n=D_n/n!$; you may use the formula
$p_n=\sum_{k=0}^n(-1)^k/k!$ from class. By convention $D_0=1$.

\topic{I. Discrete probability spaces}

\problem{A space with infinitely many outcomes}{10}
Flip a fair coin until the first head appears and record the number $K$
of flips used. Take $\Omega=\{1,2,3,\ldots\}$, let $\mathcal F=2^\Omega$,
and give outcome $k$ the mass $2^{-k}$.
\begin{enumerate}
 \item Explain where the mass $2^{-k}$ comes from, and verify that the
 masses sum to one. What does this say about the possibility that no head
 ever appears? \partpts{3}
 \item Find $\P(K\text{ is even})$ and $\P(K\text{ is divisible by }3)$.
 Find $\P(K>n)$ in two ways: by summing masses, and by describing the event
 in terms of the first $n$ flips. \partpts{4}
 \item Show that there is \emph{no} probability measure on
 $\Omega=\{1,2,3,\ldots\}$ giving every outcome the same mass. Which axiom
 does the argument use that a finite space never needs? \partpts{3}
\end{enumerate}
\begin{solution}
\begin{enumerate}
 \item $K=k$ means the flips are $T\cdots TH$ with $k-1$ tails: one
 specified string among $2^k$ equally likely strings of length $k$.
 The geometric series gives $\sum_{k\ge1}2^{-k}=\frac{1/2}{1-1/2}=1$.
 The outcomes ``first head at flip $k$'' already carry total probability
 one, so ``no head ever'' has probability zero; nothing is lost by
 omitting it from $\Omega$.
 \item Both are geometric series:
 \[
 \P(K\text{ even})=\sum_{j\ge1}4^{-j}=\frac{1/4}{1-1/4}=\frac13,
 \qquad
 \P(3\mid K)=\sum_{j\ge1}8^{-j}=\frac{1/8}{1-1/8}=\frac17 .
 \]
 Summing masses, $\P(K>n)=\sum_{k>n}2^{-k}=2^{-n}$. Directly, $K>n$
 exactly when the first $n$ flips are all tails, which has probability
 $2^{-n}$.
 \item Suppose every outcome had mass $c$. Countable additivity applied to
 $\Omega=\bigcup_k\{k\}$ gives $1=\P(\Omega)=\sum_{k\ge1}c$. If $c=0$
 the sum is $0$; if $c>0$ it diverges. Either way it is not $1$.
 The argument uses \emph{countable} additivity; on a finite space finite
 additivity suffices. (So ``a uniformly random positive integer'' is not a
 probability model.)
\end{enumerate}
\end{solution}

\problem{The larger of two dice}{10}
Roll two fair dice, one red and one green, and let $M$ be the larger of the
two values shown (the common value if they agree).
\begin{enumerate}
 \item Give $\Omega$ and the mass of each outcome. Show that
 $\P(M\le k)=k^2/36$ for $k=1,\ldots,6$, and deduce that
 $\P(M=k)=(2k-1)/36$. \partpts{4}
 \item Check that your six values of $\P(M=k)$ sum to one. Which familiar
 identity about odd numbers have you verified, and how does a $6\times6$
 grid of outcomes explain it? \partpts{3}
 \item Recording only $M$ produces a new probability space on
 $\{1,\ldots,6\}$. Find $\P(M\text{ is even})$ in that space. Compare with
 the probability that a single fair die is even, and explain the direction
 of the difference. \partpts{3}
\end{enumerate}
\begin{solution}
\begin{enumerate}
 \item $\Omega=\{1,\ldots,6\}^2$ (red, green), each of the 36 ordered pairs
 having mass $1/36$. $M\le k$ means both dice show at most $k$:
 $k^2$ outcomes. Since $\{M\le k\}$ is the disjoint union of
 $\{M\le k-1\}$ and $\{M=k\}$,
 \[
 \P(M=k)=\frac{k^2-(k-1)^2}{36}=\frac{2k-1}{36}.
 \]
 \item $(1+3+5+7+9+11)/36=36/36$. This is the identity
 $1+3+\cdots+(2n-1)=n^2$ at $n=6$. In the grid, the outcomes with $M=k$
 form an L-shaped border of $2k-1$ cells around the $(k-1)\times(k-1)$
 square, and the six borders tile the whole grid.
 \item $\P(M\text{ even})=(3+7+11)/36=21/36=7/12$, larger than $1/2$.
 The masses $(2k-1)/36$ increase with $k$, and each even value $2j$
 outweighs the odd value $2j-1$ just below it, so the even values
 collect more than half of the mass. The recorded space has six outcomes,
 but they are not equally likely.
\end{enumerate}
\end{solution}

\studentbreak
\topic{II. Permutations and combinations}

\problem{Poker hands}{12}
A hand is five cards dealt from a well-shuffled standard deck: 13 ranks
in each of 4 suits. Give each probability as an exact fraction and as a
decimal to four places.
\begin{enumerate}
 \item Describe a probability space in which a hand is an \emph{unordered}
 set of cards, and count its outcomes. \partpts{2}
 \item Find the probability of a \emph{full house}: three cards of one rank
 and two of another. \partpts{3}
 \item Find the probability of \emph{exactly one pair}: two cards of one
 rank and three cards of three other, different ranks. \partpts{4}
 \item For \emph{two pairs} (two cards of one rank, two of another, and a
 fifth card of a third rank) a student writes
 $13\binom42\cdot12\binom42\cdot44$. Explain the error, give the correct
 count, and explain why the same ``first rank, then second rank'' scheme
 was \emph{not} an error in part (b). \partpts{3}
\end{enumerate}
\begin{solution}
\begin{enumerate}
 \item $\Omega$ is the set of 5-element subsets of the deck, all equally
 likely, so $|\Omega|=\binom{52}5=2{,}598{,}960$.
 \item Choose the rank of the triple and three of its four suits, then
 the rank of the pair and two of its suits:
 $13\binom43\cdot12\binom42=3744$. The probability is
 $3744/2{,}598{,}960=6/4165\approx0.0014$.
 \item Choose the paired rank and its two suits, then three other ranks
 as an unordered set, then a suit for each:
 $13\binom42\binom{12}3 4^3=1{,}098{,}240$, giving probability
 $352/833\approx0.4226$.
 \item The scheme orders the two paired ranks, but ``kings and sevens''
 is the same hand as ``sevens and kings'': every hand is counted twice.
 The correct count is
 $\binom{13}2\binom42^2\cdot44=123{,}552$, with probability
 $198/4165\approx0.0475$. In a full house the two ranks play
 \emph{different} roles (triple and pair), so an ordered choice of ranks
 corresponds to exactly one hand. The rule: divide by the symmetry only
 when the roles are interchangeable.
\end{enumerate}
\end{solution}

\problem{Random arrangements of a word with repeats}{12}
The eleven letter tiles of $\texttt{MISSISSIPPI}$ (one M, four I's, four
S's, two P's) are placed in a uniformly random order.
\begin{enumerate}
 \item How many distinguishable arrangements are there? Explain why all
 of them are equally likely, even though the underlying experiment permutes
 eleven physical tiles. \partpts{3}
 \item Find the probability that no two S's are adjacent.
 (Arrange the other letters first, then choose gaps.) \partpts{4}
 \item Find the probability that the first and last letters are the same,
 in two ways: by counting arrangements, and by looking only at the two
 physical tiles that land in the first and last positions. \partpts{5}
\end{enumerate}
\begin{solution}
\begin{enumerate}
 \item There are $11!$ equally likely orderings of labeled tiles, and
 each distinguishable word arises from exactly $1!\,4!\,4!\,2!$ of them
 (permute identical letters among themselves). Equal-sized classes of
 equally likely outcomes are equally likely, and the number of words is
 $11!/(4!\,4!\,2!)=34{,}650$.
 \item The seven non-S letters can be arranged in $7!/(4!\,2!)=105$ ways.
 They create 8 gaps (including the two ends); the S's are non-adjacent
 exactly when they occupy four different gaps, $\binom84=70$ choices.
 Thus $105\cdot70=7350$ arrangements, and the probability is
 $7350/34{,}650=7/33$.
 \item \emph{Counting words.} If both ends are I, the middle nine letters
 are M, II, SSSS, PP: $9!/(2!\,4!\,2!)=3780$ words. Both ends S gives
 $3780$ by symmetry, and both ends P gives $9!/(4!\,4!)=630$.
 The probability is $(3780+3780+630)/34{,}650=8190/34{,}650=13/55$.

 \emph{Two tiles.} The ordered pair of tiles in positions 1 and 11 is
 uniform over the $11\cdot10$ ordered pairs of distinct tiles. The number
 of pairs with equal letters is $4\cdot3+4\cdot3+2\cdot1=26$, so the
 probability is $26/110=13/55$.
\end{enumerate}
\end{solution}

\problem{Committees and Vandermonde's identity}{10}
A five-person committee is chosen uniformly at random from a club of
6 women and 9 men.
\begin{enumerate}
 \item Find the probability that the committee has exactly two women.
 \partpts{3}
 \item Find the probability that it includes at least one woman and at
 least one man. \partpts{3}
 \item Prove \emph{Vandermonde's identity}
 $\displaystyle\sum_{j=0}^k\binom mj\binom n{k-j}=\binom{m+n}k$
 by a counting argument (with $\binom ab=0$ when $b>a$). What does it say
 about the probabilities in part (a), taken over all possible numbers of
 women? \partpts{4}
\end{enumerate}
\begin{solution}
\begin{enumerate}
 \item There are $\binom{15}5=3003$ equally likely committees, of which
 $\binom62\binom93=15\cdot84=1260$ have exactly two women. The probability
 is $1260/3003=60/143$.
 \item The complement is ``all men or all women,'' two disjoint events with
 $\binom95=126$ and $\binom65=6$ committees. The probability is
 $1-132/3003=1-4/91=87/91$.
 \item Count the $k$-person committees from $m$ women and $n$ men. Directly
 there are $\binom{m+n}k$. Alternatively classify by the number $j$ of
 women: choose them in $\binom mj$ ways and the $k-j$ men in
 $\binom n{k-j}$ ways. The classes are disjoint and exhaustive, so the
 sum equals $\binom{m+n}k$. Dividing by $\binom{m+n}k$: the probabilities
 $\binom mj\binom n{k-j}/\binom{m+n}k$ of ``exactly $j$ women,''
 $j=0,\ldots,k$, sum to one, as the probabilities of a partition of
 $\Omega$ must.
\end{enumerate}
\end{solution}

\studentbreak
\topic{III. Birthday collisions}

\problem{Two honest bounds for the birthday problem}{12}
Each of $n$ customers independently chooses a four-digit PIN uniformly from
the $d=10^4$ possibilities. Let $N_n$ be the event that all $n$ PINs are
different, so that $\P(N_n)=\prod_{k=1}^{n-1}(1-k/d)$. In class we
\emph{approximated} $\ln(1-x)\approx-x$. Here you will replace the
approximation by inequalities.
\begin{enumerate}
 \item Prove that $1-x\le e^{-x}$ for every real $x$. \partpts{3}
 \item Deduce that $\P(\text{some two PINs agree})\ge
 1-\exp\bigl(-n(n-1)/(2d)\bigr)$: the approximation from class is a
 rigorous \emph{lower} bound for the collision probability. \partpts{3}
 \item Let $E_{ij}$ be the event that customers $i$ and $j$ choose the same
 PIN. Using $\P(A\cup B)\le\P(A)+\P(B)$, show that
 $\P(\text{some two PINs agree})\le\binom n2/d$. \partpts{3}
 \item Use (c) to show that a collision is less likely than not when
 $n\le100$, and (b) to show that it is more likely than not when
 $n\ge119$ (use $\ln2\approx0.693147$). Then compute the exact
 collision probabilities for $n=118$ and $n=119$ to four decimal places.
 \partpts{3}
\end{enumerate}
\begin{solution}
\begin{enumerate}
 \item Let $g(x)=e^{-x}-(1-x)$. Then $g'(x)=1-e^{-x}$ is negative for
 $x<0$ and positive for $x>0$, so $g$ has its minimum at $x=0$, where
 $g(0)=0$. Hence $g\ge0$. (Equivalently: the convex function $e^{-x}$ lies
 above its tangent line at $0$.)
 \item Every factor satisfies $0\le1-k/d\le e^{-k/d}$, and inequalities
 between nonnegative numbers may be multiplied:
 \[
 \P(N_n)\le\prod_{k=1}^{n-1}e^{-k/d}
 =\exp\Bigl(-\frac1d\sum_{k=1}^{n-1}k\Bigr)
 =\exp\Bigl(-\frac{n(n-1)}{2d}\Bigr).
 \]
 Take complements.
 \item The collision event is $\bigcup_{i<j}E_{ij}$, a union of $\binom n2$
 events. By induction the two-event inequality gives
 $\P(\bigcup_m A_m)\le\sum_m\P(A_m)$ for finitely many events.
 Whatever PIN customer $i$ chose, customer $j$ matches it with probability
 $1/d$, so $\P(E_{ij})=1/d$ and the bound is $\binom n2/d$.
 \item For $n\le100$ the collision probability is at most
 $\binom{100}2/10^4=0.495<1/2$; it is nondecreasing in $n$, so the single
 case $n=100$ suffices. By (b), the collision probability exceeds $1/2$ as
 soon as $n(n-1)>2d\ln2\approx13{,}862.9$; since $118\cdot117=13{,}806$
 and $119\cdot118=14{,}042$, this holds for $n\ge119$. Exactly,
 \[
 1-\P(N_{118})\approx0.4999,\qquad 1-\P(N_{119})\approx0.5058,
 \]
 so the threshold is $n=119$. (To more places $1-\P(N_{118})=0.499936$,
 a remarkably near miss. The lower bound from (b) is $0.4986$ and $0.5045$
 at these two values: accurate to about one part in a thousand.)
\end{enumerate}
\end{solution}

\problem{What kind of collision?}{12}
\begin{enumerate}
 \item Four labeled people have independent birth \emph{months}, uniform on
 12 months. Classify the $12^4$ outcomes by pattern: all four different;
 exactly one shared pair (the other two different from it and from each
 other); two different shared pairs; three alike and one different; all
 four alike. Count each pattern, and check that the counts total $12^4$.
 \partpts{6}
 \item Now $n$ people have independent birthdays, uniform on $d$ days,
 with $2\le n\le d+1$. Show that the probability that exactly one pair
 shares a birthday and all other birthdays are different from it and from
 each other is $\binom n2\,d(d-1)\cdots(d-n+2)/d^n$. \partpts{3}
 \item For $n=23$ and $d=365$, compute to four decimal places the
 probabilities of no shared birthday, of the pattern in (b), and of any
 more complicated coincidence. \partpts{3}
\end{enumerate}
\begin{solution}
\begin{enumerate}
 \item In each case choose which people are grouped, then assign distinct
 months to the groups.
 \begin{center}\small
 \begin{tabular}{@{}llr@{}}
 \toprule
 Pattern & Count & \\
 \midrule
 All different & $12\cdot11\cdot10\cdot9$ & 11{,}880\\
 One pair & $\binom42\cdot12\cdot11\cdot10$ & 7{,}920\\
 Two pairs & $3\cdot12\cdot11$ & 396\\
 Three and one & $4\cdot12\cdot11$ & 528\\
 All alike & $12$ & 12\\
 \midrule
 Total & & $20{,}736=12^4$\\
 \bottomrule
 \end{tabular}
 \end{center}
 For two pairs, there are $3$ ways to split four people into two unordered
 pairs; once the split is fixed, the pair containing person 1 is
 distinguished, so its month and the other pair's month form an ordered
 choice, $12\cdot11$. (The probabilities are $55/96$, $55/144$, $11/576$,
 $11/432$, $1/1728$.)
 \item Choose the pair in $\binom n2$ ways. The pair together with the
 $n-2$ other people forms $n-1$ labeled groups needing distinct days:
 $d(d-1)\cdots(d-n+2)$ assignments. Different choices give different
 birthday lists, and all $d^n$ lists are equally likely.
 \item No match: $0.4927$. Exactly one pair: $0.3634$. Anything more
 complicated (two or more pairs, a triple, \dots):
 $1-0.4927-0.3634=0.1439$. So when a class of 23 has a coincidence, it is a
 single pair about 72\% of the time.
\end{enumerate}
\end{solution}

\studentbreak
\topic{IV. Derangements}

\problem{A recurrence for derangements}{10}
\begin{enumerate}
 \item Prove that $D_n=(n-1)\,(D_{n-1}+D_{n-2})$ for $n\ge2$, by a counting
 argument. Suggested route: in a derangement guest 1 receives some hat
 $j\ne1$; separate the cases in which guest $j$ does, or does not, receive
 hat 1. \partpts{5}
 \item Starting from $D_0=1$ and $D_1=0$, compute $D_2,\ldots,D_6$, and
 check $D_6$ against the formula for $p_6$. \partpts{2}
 \item The alternating-series estimate gives
 $|p_n-1/e|<1/(n+1)!$. Deduce that for $n\ge2$, $D_n$ is the integer
 nearest to $n!/e$, and confirm this for $n=6$. \partpts{3}
\end{enumerate}
\begin{solution}
\begin{enumerate}
 \item There are $n-1$ choices for the hat $j$ received by guest 1; by
 symmetry each contributes the same number of derangements.
 \emph{Case 1: guest $j$ receives hat 1.} Guests 1 and $j$ have swapped,
 and the remaining $n-2$ guests must derange their own $n-2$ hats:
 $D_{n-2}$ ways.
 \emph{Case 2: guest $j$ does not receive hat 1.} The guests other than 1
 must receive the hats other than $j$, with guest $j$ forbidden hat 1 and
 every other guest $i$ forbidden hat $i$. Renaming hat 1 as ``hat $j$''
 makes this exactly a derangement of $n-1$ guests: $D_{n-1}$ ways.
 The cases are disjoint and exhaustive, so
 $D_n=(n-1)(D_{n-2}+D_{n-1})$.
 \item $D_2=1\cdot(0+1)=1$, $D_3=2(1+0)=2$, $D_4=3(2+1)=9$,
 $D_5=4(9+2)=44$, $D_6=5(44+9)=265$. The formula gives
 $p_6=\frac12-\frac16+\frac1{24}-\frac1{120}+\frac1{720}
 =\frac{360-120+30-6+1}{720}=\frac{265}{720}$.
 \item Multiply by $n!$:
 $|D_n-n!/e|<n!/(n+1)!=1/(n+1)\le1/3$ for $n\ge2$. An integer within
 $1/2$ of a real number is the nearest integer to it, and it is unique.
 For $n=6$: $720/e\approx264.87$, nearest integer $265$.
\end{enumerate}
\end{solution}

\problem{How many guests get the right hat?}{12}
The hats of $n$ guests are returned as a uniformly random permutation.
\begin{enumerate}
 \item Show that the probability that \emph{exactly} $k$ guests receive
 their own hats is $\binom nk D_{n-k}/n!=p_{n-k}/k!$. \partpts{3}
 \item For $n=5$, tabulate the number of permutations with exactly $k$
 correct hats for $k=0,1,\ldots,5$, and check the total. \partpts{3}
 \item Explain in words why exactly $n-1$ correct hats is impossible, and
 how the formula in (a) knows this. For fixed $k$, find the limit of the
 probability in (a) as $n\to\infty$. \partpts{3}
 \item Show that the total number of correct hats, summed over all $n!$
 permutations, is exactly $n!$, by counting pairs (permutation, guest with
 the correct hat) in a second way. Verify this with your table from (b).
 What is the average number of correct hats, and how does it depend on
 $n$? \partpts{3}
\end{enumerate}
\begin{solution}
\begin{enumerate}
 \item Choose which $k$ guests are correct, $\binom nk$ ways; the other
 $n-k$ guests must then derange their own hats, $D_{n-k}$ ways. Different
 choices give different permutations. Dividing by $n!$,
 \[
 \frac{n!}{k!\,(n-k)!}\cdot\frac{D_{n-k}}{n!}
 =\frac1{k!}\cdot\frac{D_{n-k}}{(n-k)!}=\frac{p_{n-k}}{k!}.
 \]
 \item Using $D_0,\ldots,D_5=1,0,1,2,9,44$:
 \begin{center}\small
 \begin{tabular}{@{}l*{6}{r}r@{}}
 \toprule
 $k$ & 0 & 1 & 2 & 3 & 4 & 5 & Total\\
 \midrule
 $\binom5kD_{5-k}$ & 44 & 45 & 20 & 10 & 0 & 1 & 120\\
 \bottomrule
 \end{tabular}
 \end{center}
 \item If $n-1$ guests have their own hats, the one remaining hat belongs
 to the one remaining guest, who is then also correct. In the formula this
 is $D_1=0$. For fixed $k$, $p_{n-k}\to1/e$, so the probability tends to
 $e^{-1}/k!$. (These limits sum to $e^{-1}\sum_k1/k!=1$: they form a
 probability distribution on $\{0,1,2,\ldots\}$, which we shall meet again
 as the Poisson distribution with mean 1.)
 \item Count pairs $(\sigma,i)$ with $\sigma(i)=i$. Grouping by $\sigma$
 gives the total number of correct hats. Grouping by $i$: the permutations
 fixing $i$ are the $(n-1)!$ permutations of the other hats, so there are
 $n\cdot(n-1)!=n!$ pairs. Table check:
 $0\cdot44+1\cdot45+2\cdot20+3\cdot10+4\cdot0+5\cdot1=120$. The average
 number of correct hats is $n!/n!=1$ for every $n$: a larger party offers
 more chances, each proportionally less likely.
\end{enumerate}
\end{solution}

\par\addvspace{12pt}\Needspace{8\baselineskip}
{\bfseries\color{mint}Optional challenge: Secret Santa}\pts{10 extra}

In a Secret Santa, $n\ge3$ people draw names from a hat, and the draw is
repeated from scratch until nobody has drawn their own name.
\begin{enumerate}
 \item Explain why the final assignment is uniformly distributed over
 the $D_n$ derangements. For $n=5$, find the probability that more than
 three draws are needed. \partpts{3}
 \item Find the probability that persons 1 and 2 draw each other, in terms
 of derangement numbers, and evaluate it for $n=6$. \partpts{3}
 \item Find the probability that person 1 draws person 2, for general $n$.
 Then use (b) to find the conditional probability that 2 drew 1, given
 that 1 drew 2, for $n=6$. Is reciprocation more or less likely than the
 unconditional probability that 2 draws 1? \partpts{4}
\end{enumerate}
\begin{solution}
\begin{enumerate}
 \item Each draw is a uniform permutation, and draws are independent. The
 first successful draw is a uniform permutation conditioned on being a
 derangement; conditioning a uniform distribution on a subset gives the
 uniform distribution on that subset, since every derangement $\sigma$ has
 conditional probability $(1/n!)/(D_n/n!)=1/D_n$. For $n=5$ a draw fails
 with probability $1-44/120=19/30$, so more than three draws are needed
 with probability $(19/30)^3=6859/27{,}000\approx0.254$.
 \item If 1 and 2 swap, the other $n-2$ people must derange among
 themselves: $D_{n-2}$ of the $D_n$ equally likely derangements. For
 $n=6$ this is $D_4/D_6=9/265$.
 \item Relabeling persons $2,\ldots,n$ maps derangements to derangements,
 so person 1 is equally likely to draw each of the other $n-1$ people:
 the probability is $1/(n-1)$. For $n=6$,
 \[
 \P(2\to1\mid1\to2)=\frac{9/265}{1/5}=\frac{9}{53}\approx0.170
 <\frac15 .
 \]
 Reciprocation is slightly \emph{less} likely than the unconditional $1/5$.
 In general, by the recurrence of Problem 8 the derangements with $1\to2$
 number $D_{n-1}+D_{n-2}$, of which the $D_{n-2}$ of Case 1 reciprocate.
 The formula for $p_m$ gives $D_m=mD_{m-1}+(-1)^m$, and therefore
 \[
 \P(2\to1\mid1\to2)=\frac{D_{n-2}}{D_{n-1}+D_{n-2}}
 =\frac1{\,n+(-1)^{n-1}/D_{n-2}\,},
 \]
 which is close to $1/n$ rather than $1/(n-1)$; at $n=6$ it is
 $1/(6-\frac19)=9/53$. It is below $1/(n-1)$ for every $n\ge3$ except
 $n=4$, where the two agree.
\end{enumerate}
\textbf{Suggested credit:} as marked, (3), (3), (4); in (c) award 2 for
$1/(n-1)$ with a symmetry argument and 2 for the conditional probability
and comparison.
\end{solution}
\end{document}
